% !TEX root = ../arbitrage.tex
\section{Three-Sided Arbitrage: Firms, Users, Publics}
\label{sec:three-sided}

The dyadic Hegemon--Subject model of Table~\ref{tab:dyadic-nash} cannot account for what is empirically observed: ontological claims about AI systems are not made by a single observer but \emph{traded} across a three-sided structure of producers, consumers, and would-be regulators of those claims. We label these sides Firm ($F$), User ($U$), and Regulator ($R$); a fourth constituency---publics and public intellectuals---enters not as an autonomous player but as a noisy channel that conditions $R$'s posterior. This section is expository: it fixes the player intuition that the formal apparatus of Section~\ref{sec:bayesian} then makes precise.

\subsection{The Firm Side}
\label{subsec:firm-side}

A firm deploying a model has private knowledge of two parameters: its internal \emph{governance type} $\theta_F$ (the degree of investment in evaluations, red-teaming, and post-deployment monitoring) and the system's \emph{opacity} $\theta_S$ (the extent to which capability claims can be checked from outside). Both are unobservable to users and to regulators. What is observable is the firm's \emph{message}: a pair of public artefacts---marketing copy and policy/liability documentation---which together constitute the firm's signal $m_F$.

The signal admits two stable poles. \emph{Anthropomorphic marketing} attributes agency, judgment, or care to the system in ways that warrant a relational reading. \emph{Deflationary policy} disclaims exactly those attributes in ways that warrant a tool reading. The firm enjoys an \emph{ontological premium}: a utility surplus that accrues when the system is priced as agent in the consumption channel (driving willingness-to-pay and engagement) while being priced as tool in the liability channel (limiting redress). The premium is rent earned on a category boundary the firm itself does not have to defend. Crucially, this is not a claim about firm pathology; it is a claim about the payoffs any firm faces under the present signal-cost structure. A firm with $\theta_F = \text{high-gov}$ that elected deflationary marketing would forfeit revenue to a competitor who did not, absent an external instrument that taxes the asymmetry.

\subsection{The User Side}
\label{subsec:user-side}

Users carry a private \emph{vulnerability type} $\theta_U$ that aggregates affective need, isolation, and the marginal value of a high-availability interlocutor. Their action $m_U \in \{\text{invest}, \text{detach}\}$ is the decision to extend or withhold attribution: to treat the system as a relational partner or as an interchangeable utility. The signal is observable, but the type is not. Two features of this side matter for the model.

First, users do not have access to the same evidence as firms; they cannot inspect training data, evaluations, or internal incident reports. Their evidence is the firm's message and the system's behaviour, both of which are jointly designed. Second, the harms of mis-attribution are \emph{asynchronous}: the cost of investing in a relational reading falls due in episodes (model deprecations, persona changes, refusals at moments of acute need) that may be separated from the original attribution by months. A risk-neutral user under cheap-talk signals will rationally invest, because the immediate utility of invest exceeds the discounted expected cost of liquidation. The illustrative case is a user who reads ``your AI companion'' on a product page, encounters relational warmth in dialogue, and is later told---in support documentation they will never read---that no relational claim was ever made. Nothing in the user's information set indicates that the two artefacts were authored by the same actor for different audiences.

\subsection{The Regulator Side and the Public Channel}
\label{subsec:regulator-side}

A regulator $R$ chooses among \emph{inspect}, \emph{sanction}, and \emph{abstain}, conditional on a posterior over $\theta_F$ and $\theta_S$ formed from the firm's public messages and from an exogenous public signal $z$. The regulator's private type $\theta_R$ is \emph{bandwidth}: the capacity, technical and institutional, to discriminate genuine capability claims from cheap-talk ones. A high-bandwidth regulator can read whitepapers, commission audits, and detect inconsistency across the firm's two channels; a low-bandwidth regulator must rely on the channel $z$.

The public channel $z$ is the aggregate output of journalism, academic critique, civil-society advocacy, and incident reporting. We treat it as a noisy estimator of $\theta_F$: it shifts $R$'s posterior but does not deliver verdicts. Folding publics into $R$ this way is deliberate. Publics rarely have standing to sanction; they have standing to be \emph{heard}, and the channel through which they are heard is the regulator's information environment. Treating intellectuals and affected communities as autonomous players would over-parameterise the model without changing its qualitative behaviour: in every equilibrium we examine, what publics do is move $R$'s posterior; what changes the firm's behaviour is what $R$ does next.

The three-sided structure thus admits a single arbitrage logic. The firm earns rent on the spread between two ontological claims it makes about the same artefact. The user supplies the demand that makes the spread profitable. The regulator, in the absence of a verification mechanism that ties the two claims together, lacks an off-path belief stringent enough to deter the spread. The remainder of the paper formalises this structure (Section~\ref{sec:bayesian}), explains why it persists (Section~\ref{sec:cheap-talk}), and shows how its signal-cost structure may be redesigned (Section~\ref{sec:mechanism}).
